\bookchapter{บทที่ 33:มองอนาคต—ก้าวต่อไปของโทรศัพท์}

\chapterintro{มือถือจะไม่หายไปในชั่วข้ามคืน แต่ “สถานะผูกขาด” ของมันอาจถูกแบ่งโดยอุปกรณ์รุ่นใหม่ ในบทนี้เราจะคุยกัน 3 เรื่อง: อุปกรณ์รุ่นใหม่จะหน้าตาเป็นอย่างไร, มือถือจะ “ล่องหน” ไปเป็นรากฐานของตัวตนใหม่อย่างไร และประเด็นทางจริยธรรมที่เราต้องคิดล่วงหน้า}

\booksection{AR/VR และ Spatial Computing: “หน้าจอที่สอง” นอกเหนือจากมือถือ}

การแข่งขันของอุปกรณ์รุ่นใหม่ได้เริ่มต้นขึ้นแล้ว สิ่งที่มีแนวโน้มจะแบ่งเวลาการใช้มือถือมากที่สุดคืออุปกรณ์ AR/VR และ Spatial Computing:

\begin{itemize}
  \item \textbf{รุ่นต่อๆ ไปของ Vision Pro}:การเปิดตัว Vision Pro ของ Apple ถือเป็นจุดเริ่มต้นที่แนวคิด "Spatial Computing" ได้ลงตัวอย่างเป็นรูปธรรม แต่ราคา (เริ่มต้นที่ 29,999 หยวน) และน้ำหนักของผลิตภัณฑ์รุ่นแรกยังเป็นอุปสรรค คาดว่าในช่วงปี 2026-2027 จะมีการเปิดตัว Vision รุ่นราคาประหยัดที่เบาและถูกลง เพื่อให้ Spatial Computing ก้าวจาก "ของเล่นสำหรับนักพัฒนา" ไปสู่ "ผู้บริโภคกลุ่มแรกๆ" ซึ่งในขณะนั้น iPhone อาจกลายเป็นแหล่งพลังงานและศูนย์กลางข้อมูลให้กับแว่น Vision
  \item \textbf{แว่น AR}:เปิด "หน้าจอที่สอง" ตรงหน้าคุณ ข้อมูลดิจิทัลจะอยู่ทุกที่ราวกับอากาศหายใจ แว่นอัจฉริยะที่ Meta ร่วมมือกับ Ray-Ban ได้ขายไปแล้วหลายล้านอัน ส่วนในประเทศจีนผู้ผลิตอย่าง Rokid, XREAL, INMO ก็กำลังเปิดตัวแว่น AR ระดับผู้บริโภคเช่นกัน แว่น AR ในปี 2025-2026 สามารถแปลภาษาแบบเรียลไทม์, นำทางผ่าน HUD และคุยกับ AI Assistant ได้แล้ว แต่ยังมีข้อจำกัดเรื่องจอแสดงผล Waveguide และแบตเตอรี่ ซึ่งห่างจากการ "แทนที่มือถือ" อีก 3-5 ปี
  \item \textbf{แหวนอัจฉริยะและอุปกรณ์สวมใส่}:แหวนอัจฉริยะอย่าง Samsung Galaxy Ring และ Oura Ring เริ่มได้รับความนิยม ทำหน้าที่เป็น "เซ็นเซอร์ติดตัว" ที่เสริมความสามารถในการตรวจสอบสุขภาพของมือถือ บทบาทของมือถือกำลังเปลี่ยนจาก "อุปกรณ์รอบรอบด้าน" ไปเป็นการผสมผสานระหว่าง "อุปกรณ์หลัก + อุปกรณ์เสริมขนาดเล็กหลายชิ้น"
\end{itemize}

\booksection{Brain-Computer Interface: ความคิดกลายเป็นสื่อกลางในการโต้ตอบ}

ความคิดจะกลายเป็นสื่อกลางในการโต้ตอบด้วยตัวมันเอง แต่เส้นทางนี้ยาวนานกว่าที่คิด:

\begin{itemize}
  \item \textbf{ความคืบหน้าล่าสุดของ Neuralink}:Neuralink ของ Elon Musk ได้ทำการฝังในร่างกายมนุษย์ครั้งแรกในปี 2024 ผู้รับการทดลองคนแรก Noland Arbaugh สามารถควบคุมเมาส์คอมพิวเตอร์ด้วยความคิดเพื่อเล่นหมากรุกสากล ปัจจุบัน BCI (Brain-Computer Interface) แบบฝังในร่างกายมุ่งเน้นไปที่การช่วยเหลือผู้ป่วยอัมพาตเป็นหลัก โดยมีเป้าหมายเพื่อช่วยให้พวกเขาสามารถ "ควบคุม" โลกดิจิทัลได้อีกครั้ง
  \item \textbf{BCI แบบไม่ผ่าตัด}:หลายบริษัทกำลังพัฒนาห่วงศีรษะแบบไม่ผ่าตัด (เช่น Kernel, OpenBCI) ที่สามารถอ่านสัญญาณคลื่นสมองได้โดยไม่ต้องผ่าตัด แต่ความแม่นยำยังด้อยกว่าแบบฝังในร่างกายมาก ซึ่งยังห่างไกลจากการ "พิมพ์ด้วยความคิด"
  \item \textbf{มุมมองระดับผู้บริโภค}:Brain-Computer Interface ยังห่างจากการโต้ตอบกับมือถือระดับผู้บริโภคอีก 10-20 ปี แต่ในระยะสั้นอาจถูกนำไปใช้ในงานช่วยเหลือผู้พิการ หรือการตรวจสอบระดับสมาธิก่อน ประเด็นด้านจริยธรรมและความเป็นส่วนตัว ("คลื่นสมองของฉันถือเป็นความลับของใคร") จะกลายเป็นประเด็นถกเถียงก่อนที่เทคโนโลยีจะพัฒนาจนสมบูรณ์
\end{itemize}

\booksection{การ “มองไม่เห็น” ของสมาร์ทโฟนและรากฐานของอัตลักษณ์ดิจิทัลใหม่}

อีกหนึ่งทิศทางการวิวัฒนาการในอนาคตของสมาร์ทโฟนอาจเป็นการ “มองไม่เห็น” — กลายเป็น “ผู้อยู่เบื้องหลัง” ที่อยู่ทุกหนทุกแห่ง และเป็นรากฐานเดียวของอัตลักษณ์ดิจิทัลของคุณ:

\begin{itemize}
  \item \textbf{การยืนยันสองขั้นผ่านคลาวด์}:ยืนยันตัวตน "คุณคือคุณ" ได้ในพริบตา เข้าสู่ระบบบริการใด ๆ ก็ตามโดยไม่ต้องกรอกรหัสยืนยันด้วยมือ
  \item \textbf{มรดกดิจิทัลกลายเป็นค่าเริ่มต้นจากโรงงาน}:ตั้งค่าแผนสืบทอดสินทรัพย์ดิจิทัลไว้ล่วงหน้าตั้งแต่เปิดเครื่องใหม่ ทำให้ "เรื่องหลังจากจากไป" กลายเป็นค่าเริ่มต้นจากโรงงาน
  \item \textbf{ที่2เส้น:ฉันของบุคคลข้อมูลใช่ฉันของการเงินผลิต คุณของหัวต้องการตำแหน่งรับผิดชอบใช่ปลอดภัยเข้ารหัสและสะดวกรวดเร็วซิงค์}:ตู้นิรภัยข้อมูลส่วนบุคคลบน Cloud ของคุณต่างหากที่เป็นศูนย์กลาง ส่วนฮาร์ดแวร์เป็นเพียงช่องทางเข้าถึง เปลี่ยนเครื่อง เปลี่ยนอุปกรณ์ ตัวตนและข้อมูลจะตามคุณไปเสมอ
\end{itemize}

\booksection{การขบคิดด้านจริยธรรม: เรา หรือสมาร์ทโฟน กันแน่ที่เป็นเจ้าของ?}

คำถามสุดท้ายที่หลอกหลอนไม่หายก็คือ: เราเป็น “เทพ” ที่เทคโนโลยีเสริมพลัง หรือเป็น “สมองในถัง” ที่อัลกอริทึมเลี้ยงดู?

\begin{itemize}
  \item \textbf{จริยธรรม AI}:เมื่อ AI บนอุปกรณ์เข้าใจความชอบของคุณดีกว่าตัวคุณเอง "การเลือก" ของคุณเป็นเจตจำนงอิสระ หรือเป็นผลลัพธ์ที่ถูกชี้นำ? เมื่อ AI Agent สั่งซื้อสินค้าหรือตอบกลับข้อความแทนคุณ ความรับผิดชอบและขอบเขตอยู่ที่ไหน?
  \item \textbf{อธิปไตยของข้อมูล}:ข้อมูลสุขภาพ ข้อมูลคลื่นสมอง ข้อมูลลักษณะทางชีวภาพของคุณ เป็นของคุณ ของผู้ผลิตอุปกรณ์ หรือของผู้ให้บริการ Cloud? พระราชบัญญัติคุ้มครองข้อมูลส่วนบุคคลและ GDPR ของสหภาพยุโรปเป็นเพียงจุดเริ่มต้น การนำ "สิทธิในการเคลื่อนย้ายข้อมูล" และ "สิทธิในการถูกลืม" มาใช้จริงยังมีอีกยาวไกล
  \item \textbf{ช่องว่างดิจิทัล}:สมาร์ทโฟน AI, หน้าจอพับได้, การสื่อสารผ่านดาวเทียม ล้วนกำลังขยายช่องว่างระหว่าง "ใช้ได้" และ "ใช้ไม่ได้" การเข้าถึงโครงสร้างพื้นฐานดิจิทัลอย่างเป็นธรรมจะกลายเป็นประเด็นทางสังคมใหม่ — สำคัญเหมือนกับเมื่อ 20 ปีก่อนที่ว่า "รู้จักใช้อินเทอร์เน็ตหรือไม่"
  \item เมื่ออัลกอริทึมรู้จักความชอบของคุณดีกว่าตัวคุณ "ถูกใจ" ของคุณมีสักเท่าไรที่มาจากใจจริง?
  \item เมื่อคุณไม่รู้ตัวจุดไฟหน้าจอทุกไม่กี่นาที นั่นคือนิสัยของคุณ หรือ "พฤติกรรมติดใจ" ที่ถูกออกแบบมา?
\end{itemize}

\booksection{เชื่อมต่อแต่ยังเป็นอิสระ — ปรัชญาการใช้ชีวิตดิจิทัลของคุณ}

ไม่ว่าเทคโนโลยีจะวิวัฒนาการไปเช่นไร ปรัชญาการใช้ชีวิตดิจิทัลต่อไปนี้ แม้อีก 10 ปีก็ยังใช้ได้:

\begin{enumerate}
  \item : เป็นผู้จัดสรรทรัพยากร ไม่ใช่ผู้รับสิ่งที่ถูกผลักดันมาให้
  \item ให้ระบบอัตโนมัติเป็น "ผู้นั่งเบียะ" ที่ซื่อสัตย์ของคุณ
  \item สร้าง "แนวหลังดิจิทัล" ที่มั่นคงของคุณ
  \item ลงมือ "ปลดระวางของดิจิทัล" และ "ถือศีลงดเว้นสังคม" เป็นระยะ
  \item ยึดความปลอดภัยและความเป็นส่วนตัวเป็นเงื่อนไขก่อนของความสะดวกทุกอย่างตลอดไป
\end{enumerate}

\begin{tcbpractice}{ลองดูสิ: ร่าง "ข้อกำหนดการใช้อุปกรณ์ดิจิทัลส่วนบุคคล" ของคุณ}
เขียน “ข้อกำหนดการให้บริการ” ง่าย ๆ สำหรับคุณและสมาร์ทโฟน (รวมถึงแว่น AR และผู้ช่วย AI ในอนาคต) เพื่อกำหนดความสัมพันธ์ระหว่างคุณ ตัวอย่างเช่น:

\begin{itemize}
  \item : ข้อ 1: เวลาความสนใจของฉันเป็นของฉัน เวลาฉันทำงานหรืออยู่กับครอบครัว คุณต้องเข้าสู่ "โหมดห้ามรบกวน" เองโดยอัตโนมัติ
  \item : ข้อ 2: ข้อมูลส่วนตัวของฉันเป็นทรัพย์สินของฉัน หน้าที่แรกของคุณคือการเข้ารหัสอย่างปลอดภัยและซิงก์อย่างสะดวก
  \item : ข้อ 3: เป้าหมายร่วมกันของเราคือทำให้ชีวิตมีประสิทธิภาพและประสบการณ์ดีขึ้น ถ้าคุณใช้ดีไซน์ชวนหลงจนทำให้ฉันติดหนึบ ฉันมีสิทธิ์ถอนการติดตั้งคุณ
\end{itemize}

\end{tcbpractice}

\begin{tcbsummary}{สรุปบทนี้}
\begin{itemize}
  \item Vision Pro และแว่น AR เป็นช่องทางเข้าถึงการประมวลผลใหม่นอกเหนือจากสมาร์ทโฟน แต่ในระยะสั้นยังไม่สามารถแทนที่สมาร์ทโฟนได้
  \item Brain-Computer Interface ยังห่างจะระดับผู้บริโภคอีก 10-20 ปี ปัญหาด้านจริยธรรมมาก่อนที่เทคโนโลยีจะพัฒนาสุกงอม
  \item สมาร์ทโฟนกำลังวิวัฒนาการไปสู่ทิศทาง "มองไม่เห็น" จากการมีอุปกรณ์เป็นศูนย์กลางสู่การมีตัวตนเป็นศูนย์กลาง
  \item จริยธรรม AI, อธิปไตยของข้อมูล, ช่องว่างดิจิทัล คือประเด็นหลักใน 10 ปีข้างหน้า
  \item เชื่อมต่อแต่ยังเป็นอิสระ คือปรัชญาที่ไม่เปลี่ยนแปลงของการใช้ชีวิตดิจิทัล
\end{itemize}

\end{tcbsummary}
